\chapter{Pruebas y evaluación}

\section{Enfoque de evaluación}

La validación de este trabajo no podía limitarse a comprobar que el código se
ejecuta sin errores. Esa parte queda cubierta por las pruebas automatizadas,
pero no basta para saber si el laboratorio cumple su objetivo principal:
ayudar a estudiar paradigmas de toma de decisiones humanas a partir de los
modelos generados en la primera fase del proyecto.

El laboratorio combina dos naturalezas muy distintas que no se pueden medir
con la misma vara. Por un lado, un \textbf{núcleo determinista} ---el motor de
simulación, el Tracker y la persistencia en el Knowledge Backbone--- cuyo
comportamiento es reproducible y se puede comprobar contra una verdad
programática. Por otro, una \textbf{capa de juicio basada en modelos de
lenguaje} ---Architect, Analyst y Reporter--- cuyas salidas son abiertas y no
tienen una única respuesta correcta. Evaluar ambas igual llevaría o bien a
exigir determinismo donde no lo hay, o bien a aceptar como válido cualquier
texto plausible.

Por eso la evaluación se organiza en \textbf{tres capas}, siguiendo la
taxonomía de métodos de cómputo de métricas para agentes basados en modelos de
lenguaje de Mohammadi et al.~\cite{mohammadi2025agenteval}, que distingue tres
familias complementarias: comprobaciones basadas en código (las más
deterministas y objetivas), juicio mediante modelo de lenguaje
(\textit{LLM-as-a-judge}) y revisión humana (el patrón oro para lo subjetivo).
Esta separación es coherente con la práctica clásica de verificación y
validación de modelos de simulación~\cite{sargent2013vv}: la verificación
comprueba que el sistema está bien construido; la validación, que produce
resultados representativos del fenómeno estudiado.

\begin{table}[H]
\begin{center}
\small
\begin{tabular}{|>{\raggedright\arraybackslash}p{2.7cm}|>{\raggedright\arraybackslash}p{5.2cm}|>{\raggedright\arraybackslash}p{4.4cm}|} \hline
\textbf{Capa} & \textbf{Qué valida} & \textbf{Agentes / componentes} \\ \hline
1. Núcleo determinista & Contratos técnicos contra verdad programática (métricas duras, binarias) & Motor de simulación, Tracker, persistencia \\ \hline
2. Comportamiento (\textit{golden scenarios}) & Que los modelos se comportan como predice la teoría que implementan & Modelos de fase 1, Analyst \\ \hline
3. Juicio y revisión humana & Calidad de las salidas abiertas frente a criterio experto & Architect, Analyst, Reporter \\ \hline
\end{tabular}
\caption{Marco de evaluación por capas y su correspondencia con los componentes del laboratorio.}
\label{tab:eval-capas}
\end{center}
\end{table}

El mensaje que articula el capítulo es que la primera fase valida que el
\emph{modelo funciona} ---compila, pasa tests y se integra---, mientras que la
segunda fase valida que el \emph{modelo se comporta como predice la teoría} y
que la infraestructura lo \emph{observa con fidelidad}. Esa es precisamente la
brecha que la fase 1 deja abierta: que un modelo se ejecute sin errores no
garantiza que reproduzca el fenómeno que pretende
representar~\cite{wilson2019tenrules}.

\section{Capa 1: núcleo determinista}

Esta capa reúne las comprobaciones objetivas y reproducibles. La literatura de
evaluación de agentes recomienda preferir comprobaciones deterministas siempre
que exista una verdad programática, porque los jueces basados en modelos de
lenguaje son no deterministas y tienden a premiar rasgos estilísticos por
encima de la corrección~\cite{mohammadi2025agenteval}. Por eso el núcleo
determinista es la primera capa, y el juicio mediante modelo de lenguaje se
reserva para las salidas genuinamente abiertas.

\subsection{Contrato observable de los modelos}

Antes de medir comportamiento hay que garantizar que los modelos respetan el
contrato común en tiempo de ejecución, no solo en sus pruebas unitarias. Se
comprueban tres invariantes sobre cada modelo cargado dinámicamente:

\begin{enumerate}
\item \textbf{\texttt{decide} es de solo lectura.} Invocar \texttt{decide} dos
veces sobre la misma percepción no modifica el estado interno: el
\texttt{get\_state} previo y posterior son idénticos. Esto evita que el modelo
aprenda o cambie variables antes de que el entorno aplique la acción.
\item \textbf{\texttt{update} es el único mutador.} El estado solo cambia tras
recibir recompensa y nueva percepción.
\item \textbf{\texttt{get\_state} expone \texttt{q\_values}.} La presencia de
esta clave es lo que permite a la capa de observación representar alternativas
de acción y confianza relativa de forma uniforme para cualquier modelo.
\end{enumerate}

Esta comprobación es binaria y no requiere juicio. Corresponde a las reglas
mínimas de sanidad previas a cualquier ajuste de modelo de
comportamiento~\cite{wilson2019tenrules}, trasladadas al contrato de ejecución
del laboratorio.

\subsection{Determinismo del motor}

El determinismo del motor sostiene la comparativa multi-modelo: con la misma
semilla y la misma configuración, dos simulaciones producen eventos idénticos.
La parte determinista del sistema es el motor; la variabilidad de los modelos
de lenguaje se discute en las limitaciones. Esta propiedad se verifica de forma
automática comparando dos ejecuciones byte a byte sobre la secuencia de eventos
emitidos.

\subsection{Fidelidad de la persistencia: la tripleta \textit{joinable}}

Cada observación de simulación debe materializarse de forma consistente en los
tres almacenes del Knowledge Backbone: una fila en PostgreSQL, un vector denso
en Qdrant y un vector \textit{sparse} en Qdrant, siempre bajo el mismo
\textit{UUID}. La fidelidad del Tracker se mide como la igualdad entre el número
de observaciones generadas y el incremento que produce la ejecución en cada uno
de los tres almacenes.

\begin{table}[H]
\begin{center}
\small
\begin{tabular}{|>{\raggedright\arraybackslash}p{7.2cm}|>{\raggedright\arraybackslash}p{4.8cm}|} \hline
\textbf{Métrica de persistencia} & \textbf{Valor} \\ \hline
Observaciones de simulación generadas & $N$ \\ \hline
Filas nuevas en \texttt{simulation\_observations} (PostgreSQL) & $N$ \\ \hline
Memorias densas escritas (Voyage, 1024 dim.) & $N$ \\ \hline
Memorias \textit{sparse} escritas (BM25 nativo) & $N$ \\ \hline
\end{tabular}
\caption{Fidelidad de la persistencia: la igualdad $N=N=N$ valida la tripleta \textit{joinable} (ninguno de los tres almacenes queda desincronizado).}
\label{tab:joinable}
\end{center}
\end{table}

La igualdad de las tres últimas filas, medida como el incremento que produce la
ejecución sobre el estado previo del backbone, valida la \textbf{tripleta
joinable} descrita en el diseño: la persistencia no deja ninguno de los tres
almacenes desincronizado.

\subsection{Suite de regresión}

La suite automática se usa como red de regresión para asegurar que la
evaluación de las capas superiores se apoya en un sistema estable. Combina
\texttt{uv run pytest} para dominio, agentes, Orchestrator, persistencia,
recuperación y carga de modelos; Playwright para la sesión completa del cliente
web y las vistas del Knowledge Backbone; y compilación LaTeX para la propia
memoria. Los cambios visibles en la interfaz se revisaron además manualmente
sobre la aplicación viva, porque una prueba \textit{end-to-end} puede confirmar
que un elemento existe sin juzgar su legibilidad.

\begin{table}[H]
\begin{center}
\small
\begin{tabular}{|>{\raggedright\arraybackslash}p{4.6cm}|>{\raggedright\arraybackslash}p{1.6cm}|>{\raggedright\arraybackslash}p{1.4cm}|>{\raggedright\arraybackslash}p{4.0cm}|} \hline
\textbf{Bloque} & \textbf{Estado} & \textbf{Casos} & \textbf{Qué demuestra} \\ \hline
Backend (\textit{pytest}, 36 archivos) & PASS & 326/326 & Dominio, agentes, Orchestrator, persistencia, recuperación y carga de modelos. \\ \hline
Determinismo del motor (misma semilla) & PASS & 1/1 & Dos ejecuciones con la misma semilla y configuración producen eventos idénticos. \\ \hline
E2E web (modo \textit{mock}) & PASS & 4/4 & Pipeline completo en la interfaz (spec, simulación, \textit{replay}, análisis, informe) sin servicios externos. \\ \hline
E2E web (sesión real y backbone) & Separada & 13 & Sesión completa y vistas del Knowledge Backbone; requieren el \textit{backend} en ejecución. \\ \hline
\end{tabular}
\caption{Resultado de la ejecución de la suite automatizada (red de regresión de la capa 1).}
\label{tab:suite}
\end{center}
\end{table}

Esta batería importa dinámicamente los modelos producidos por el Builder y
comprueba que cumplen el \texttt{Protocol} declarado en
\texttt{simlab.environment} antes de iniciar una simulación. No compara
rendimiento con otros sistemas; sirve como evidencia interna de que los
contratos técnicos del laboratorio se mantienen estables.

\section{Capa 2: evaluación de comportamiento (\textit{golden scenarios})}

La capa 1 demuestra que el sistema es técnicamente correcto, pero no que los
modelos hagan lo que la teoría predice. Para cerrar esa brecha se emplea la
técnica de los \textit{golden scenarios}: escenarios controlados en los que
\emph{se conoce de antemano la verdad}, porque viene fijada por la teoría que el
modelo dice implementar. Se mide entonces si el laboratorio ---y en particular
el Analyst--- recupera esa verdad a partir de la simulación. Este enfoque es el
recomendado para evaluar agentes cuando no existe una única respuesta
correcta~\cite{mohammadi2025agenteval}, y traslada al laboratorio la práctica de
validación de modelos de comportamiento mediante recuperación de patrones
conocidos~\cite{wilson2019tenrules} y
falsación~\cite{palminteri2017falsification}.

\subsection{Modelos sujeto y \textit{baselines}}

Los sujetos de evaluación son los cuatro modelos generados por la primera fase
en la ejecución de referencia, repartidos en dos paradigmas con predicciones
cuantitativas falsables:

\begin{table}[H]
\begin{center}
\small
\begin{tabular}{|>{\raggedright\arraybackslash}p{3.0cm}|>{\raggedright\arraybackslash}p{4.4cm}|>{\raggedright\arraybackslash}p{4.6cm}|} \hline
\textbf{Paradigma} & \textbf{Modelo} & \textbf{Predicción teórica} \\ \hline
Forrajeo óptimo & Teorema del valor marginal (MVT) & Abandona el parche cuando la tasa marginal cae bajo la media del entorno~\cite{charnov1976mvt} \\ \hline
Forrajeo óptimo & Amplitud de dieta (\textit{diet breadth}) & Regla cero-uno: ignora la presa de baja rentabilidad si abunda la buena~\cite{macarthur1966dietbreadth} \\ \hline
Aprendizaje por refuerzo & Q-learning tabular ($\varepsilon$-greedy) & Curva de aprendizaje: el rendimiento mejora con la experiencia \\ \hline
Aprendizaje por refuerzo & \textit{Actor-critic} (softmax) & Aprende; el error TD actúa como señal de recompensa-predicción \\ \hline
\end{tabular}
\caption{Modelos sujeto de los \textit{golden scenarios} y la predicción falsable que cada uno debe reproducir.}
\label{tab:modelos-sujeto}
\end{center}
\end{table}

Para anclar las métricas de comportamiento se añaden dos modelos de referencia
de verdad conocida, implementados a mano en el paquete \texttt{benchmark}:
\texttt{GreedyForagerOracle}, que se dirige siempre a la comida más cercana y la
consume (define el \textbf{techo} de eficiencia de forrajeo), y
\texttt{RandomModel}, que elige una acción uniformemente al azar (define el
\textbf{suelo}). Ambos respetan el contrato común por \textit{duck typing}. Un
modelo generado que forrajee de verdad debe situarse cerca del oráculo y, en
todo caso, batir claramente al azar; si no bate al azar, su comportamiento es
indistinguible de la casualidad y el escenario ha fallado en detectar un patrón
real.

\subsection{Escenarios}

\begin{table}[H]
\begin{center}
\small
\begin{tabular}{|>{\raggedright\arraybackslash}p{1.6cm}|>{\raggedright\arraybackslash}p{4.8cm}|>{\raggedright\arraybackslash}p{5.4cm}|} \hline
\textbf{ID} & \textbf{Predicción falsable} & \textbf{Métrica observada} \\ \hline
GS-OFT-1 & Abandono de parche al agotarse & Ciclos residencia$\rightarrow$partida; \texttt{patch\_residence\_time} crece y se reinicia \\ \hline
GS-OFT-2 & $\uparrow$ coste de viaje $\Rightarrow$ $\uparrow$ residencia & Tiempo medio de residencia al partir, con coste de viaje bajo frente a alto \\ \hline
GS-OFT-3 & Regla cero-uno de inclusión de dieta & Discontinuidad en la inclusión de presa de baja palatabilidad según abundancia de la buena \\ \hline
GS-RL-1 & Curva de aprendizaje (RL aprende; OFT es plano) & Recompensa por ventana temporal; \texttt{td\_error}$\rightarrow 0$; \texttt{exploration\_rate} decae \\ \hline
\end{tabular}
\caption{\textit{Golden scenarios} y su predicción teórica de verdad conocida.}
\label{tab:golden-scenarios}
\end{center}
\end{table}

El escenario más discriminativo es \textbf{GS-OFT-2}, la estática comparativa
central del teorema del valor marginal~\cite{charnov1976mvt}: al aumentar el
coste de desplazarse entre parches, un forrajeador óptimo debe permanecer más
tiempo en cada parche antes de abandonarlo. Es una predicción direccional cuyo
signo se conoce por la teoría, de modo que se verifica si el modelo lo
reproduce ---en el espíritu de la recuperación de
parámetros~\cite{wilson2019tenrules}. \textbf{GS-RL-1} es discriminativo entre
paradigmas: los modelos de aprendizaje por refuerzo deben mejorar con la
experiencia, mientras que los de forrajeo óptimo, que no aprenden, deben
mostrar un rendimiento plano. El contraste entre ambos comportamientos es, en sí
mismo, una prueba de falsación~\cite{palminteri2017falsification}.

\textbf{TODO:} ejecutar la batería de \textit{golden scenarios} sobre los cuatro
modelos y los dos \textit{baselines}, y recoger en una tabla los valores
medidos junto con el resultado esperado por la teoría (PASS/FALLA por
predicción). Estos resultados son reproducibles a partir del estado de la
primera fase restaurado en local.

\subsection{Recuperación de la verdad por el Analyst}

El test propio de esta capa para la parte de modelo de lenguaje es: dado cada
escenario, ¿atribuye el Analyst el comportamiento al paradigma correcto
(forrajeo óptimo frente a aprendizaje) y señala el patrón que la teoría predice?
Como la verdad de cada escenario es conocida, la calidad del Analyst se puntúa
con precisión y exhaustividad (\textit{precision} y \textit{recall}) frente a esa
verdad, y no con una valoración subjetiva. El Knowledge Backbone refuerza este
contraste: la run restaurada contiene el postulado del teorema del valor
marginal vinculado a su evidencia empírica (Charnov, 1976), de modo que el
Analyst puede contrastar el patrón observado contra el postulado recuperado.

\textbf{TODO:} definir el conjunto de etiquetas de verdad por escenario y
reportar \textit{precision}/\textit{recall} de las atribuciones del Analyst.

\subsection{Validación de los modelos generados}

Más allá de evaluar al Analyst, los \textit{golden scenarios} permiten validar
la \emph{ciencia} de los modelos generados, no solo su ejecución. Se siguen las
reglas de validación de modelos de comportamiento de Wilson y
Collins~\cite{wilson2019tenrules}: comprobar que el modelo reproduce, en
simulación, los patrones cualitativos que la teoría predice (un análogo de la
comprobación predictiva \textit{a posteriori}), y que dos paradigmas distintos
producen comportamientos distinguibles (un análogo de la recuperación de
modelo). La importancia de que un modelo pueda además \emph{fallar} donde otro
acierta ---la falsación--- justifica el uso de los \textit{baselines} como
controles negativos~\cite{palminteri2017falsification}.

\section{Capa 3: juicio mediante modelo de lenguaje y revisión humana}

Las salidas del Architect, el Analyst y el Reporter son abiertas: no hay una
única especificación de entorno correcta, ni un único análisis válido, ni un
único informe. Para evaluarlas se combina el juicio mediante modelo de lenguaje
con la revisión humana como patrón oro.

\subsection{\textit{LLM-as-a-judge}}

El uso de un modelo de lenguaje como juez es un método validado: jueces fuertes
alcanzan más de un 80\,\% de acuerdo con las preferencias humanas, al mismo
nivel que el acuerdo entre humanos~\cite{zheng2023judge}. Ahora bien, ese
acuerdo es específico de cada tarea y dominio, y los jueces presentan una
taxonomía de sesgos bien documentada ---posición, verbosidad,
\textit{self-enhancement}, entre otros--- que amenaza su
validez~\cite{calm2024biases}. Para mitigarlos se adoptan tres prácticas: usar
\textbf{un juez aislado por dimensión} evaluada en lugar de uno único para todo,
\textbf{calibrar} el juez contra criterio experto antes de confiar en él, y
\textbf{anclar el juicio en rúbricas estructuradas} explícitas.

\subsection{Meta-evaluación: acuerdo juez--humano}

La fiabilidad del propio juez se mide contrastando sus veredictos con los
humanos, considerados el patrón oro, mediante métricas de acuerdo y correlación
---exactitud, $\kappa$ de Cohen, $\rho$ de Spearman--- según el tipo de
dato~\cite{li2024judges}. Solo cuando ese acuerdo es suficientemente alto se
reduce la proporción de revisión humana.

\textbf{TODO:} fijar el conjunto de salidas etiquetadas por el tutor, calcular
el acuerdo juez--humano y reportar el umbral a partir del cual se confía en el
juez.

\subsection{Fidelidad del Reporter}

El Reporter genera informes en PDF a partir de los datos de la simulación y del
análisis. El riesgo específico es la \emph{alucinación}: que el informe afirme
números o hechos que no se corresponden con los datos. Para acotarlo se mide la
\textbf{consistencia factual} del informe frente a su fuente, con técnicas de
verificación basadas en preguntas y respuestas ---comprobar que las preguntas
hechas sobre el informe y sobre los datos producen la misma
respuesta~\cite{wang2020qags}--- o de inferencia textual a nivel de
afirmación~\cite{laban2022summac}.

\textbf{TODO:} seleccionar las magnitudes clave de cada informe (energía final,
comida consumida, número de episodios) y comprobar automáticamente que coinciden
con los datos de simulación.

\subsection{Revisión experta}

La capa humana es el patrón oro contra el que se calibran las anteriores. La
evaluación final no es puramente cuantitativa: el criterio decisivo es si el
tutor y el grupo de investigación consideran que la herramienta les reduce
fricción para explorar paradigmas nuevos. Por eso la revisión cualitativa forma
parte de la validación: permite detectar si el sistema produce información útil o
si solo presenta una ejecución técnicamente correcta. En cada escenario se
recoge la misma clase de evidencia ---configuración del entorno, modelos usados,
\textit{replay} de la simulación, eventos críticos, salida del Tracker, análisis
del Analyst, informe del Reporter y comentarios del grupo--- de modo que cada
ejecución pueda revisarse después y no quede como una interacción efímera en el
navegador.

\section{Escenario principal de validación \textit{end-to-end}}

Por encima de las capas anteriores, un escenario principal valida el recorrido
completo como herramienta. Compara dos modelos cargados desde la primera fase
que pertenecen a paradigmas distintos: el del \textbf{teorema del valor marginal}
(forrajeo óptimo) y el de \textbf{Q-learning tabular} (aprendizaje por refuerzo).
La elección no es arbitraria: combina un objetivo interpretable, variables
internas observables ---tiempo de residencia y tasa marginal en uno, error TD y
tasa de exploración en el otro--- y diferencias esperables entre modelos. Si el
sistema funciona, no debería limitarse a mostrar que ambos agentes se mueven por
la rejilla, sino permitir comparar cómo cada modelo responde a la disponibilidad
de recursos y cómo evoluciona su estado interno.

La ejecución recorre el pipeline completo:

\begin{enumerate}
\item el Orchestrator inicia la sesión y lista los modelos disponibles;
\item el Architect entrega una especificación de entorno coherente con el
paradigma de forrajeo;
\item el motor de simulación ejecuta los modelos y publica eventos en vivo;
\item el Tracker produce trayectorias y episodios relevantes;
\item el Analyst resume patrones, trazas de decisión y gráficas de energía,
acciones y valores internos;
\item el Reporter genera un informe PDF descargable para la ejecución.
\end{enumerate}

\subsection{Criterios de evaluación}

Para que el escenario se considere satisfactorio, la ejecución debe cumplir
cinco criterios.

\begin{enumerate}
\item \textbf{Continuidad del flujo.} El usuario debe poder pasar de una
petición conversacional a una simulación observable, después a un análisis y
finalmente a un informe descargable, sin cambiar de herramienta.
\item \textbf{Integración con la primera fase.} Los modelos generados por el
Builder deben cargarse dinámicamente y ejecutarse mediante el contrato común
\texttt{decide}, \texttt{update} y \texttt{get\_state}, sin adaptadores
específicos para cada modelo.
\item \textbf{Observabilidad de la simulación.} La interfaz debe mostrar el
estado del episodio, los eventos emitidos por el backend, las trayectorias de
los agentes y las variables internas relevantes.
\item \textbf{Utilidad del análisis.} La salida del Analyst debe ayudar a
entender diferencias entre modelos, no solo repetir datos de la simulación. En
particular, debe señalar patrones de comportamiento, episodios críticos y
relaciones entre acciones, recompensas y estado interno.
\item \textbf{Trazabilidad del resultado.} El informe generado debe recoger el
contexto de la simulación, los resultados principales y la interpretación del
análisis de forma revisable.
\end{enumerate}

Estos criterios cierran el recorrido iniciado en requisitos. La continuidad del
flujo cubre el chat y la interfaz reactiva (RF-10, RF-12 y RNF-04); la
integración con la primera fase valida la carga dinámica y el desacoplamiento
(RF-01, RNF-01 y RNF-02); la observabilidad de la simulación cubre el registro
estructurado y la trazabilidad operativa (RF-04, RNF-03 y RNF-09); la utilidad
del análisis recoge el contraste con el Knowledge Backbone y las visualizaciones
(RF-05, RF-06 y RF-08); y la trazabilidad del resultado valida el informe y sus
artefactos asociados (RF-07 y RNF-03).

\subsection{Métricas de una ejecución real}

Se instrumentó una ejecución real del escenario contra un modelo de lenguaje
servido a través de OpenRouter y sobre los servicios reales del Knowledge
Backbone, recorriendo el pipeline completo ---Architect, motor de simulación,
Tracker, Analyst y Reporter.

\textbf{TODO:} reinstrumentar esta ejecución con los modelos del teorema del
valor marginal y Q-learning, y reportar latencia de extremo a extremo,
configuración del entorno generado por el Architect, número de observaciones
persistidas y la igualdad de la tripleta \textit{joinable}
(tabla~\ref{tab:joinable}). Para la defensa, la ejecución en directo usará un
\textit{backend} con los servicios externos sustituidos por respuestas
deterministas, de forma que la demostración no dependa de la latencia ni de la
disponibilidad de proveedores de modelos de lenguaje.

\section{Limitaciones}

La suite y los escenarios descritos validan propiedades técnicas y de
comportamiento, pero conviene delimitar su alcance.

La primera limitación es la diferencia entre validación estructural y validación
semántica. Las pruebas de la capa 1 comprueban que el Analyst produce un JSON
con la forma esperada y que el Architect genera una especificación válida contra
el modelo formal del entorno, pero no que el patrón identificado sea la mejor
interpretación de las trayectorias ni que la especificación sea la más adecuada
para el paradigma descrito. Las capas 2 y 3 acotan esa brecha ---los
\textit{golden scenarios} con verdad conocida y la calibración del juez contra
criterio experto---, pero la adecuación psicológica o neurocientífica última de
un modelo sigue requiriendo revisión experta del tutor y del grupo de
investigación~\cite{wilson2019tenrules}.

La segunda limitación es la reproducibilidad acotada del comportamiento de los
modelos de lenguaje. La idempotencia se garantiza para el motor de simulación,
siempre que se use la misma semilla y la misma configuración, pero las salidas
generadas por modelos de lenguaje mantienen cierta variabilidad. Esta es la
razón de fondo por la que el juicio de la capa 3 nunca sustituye del todo a la
revisión humana, sino que se calibra contra ella~\cite{zheng2023judge}.

La tercera pertenece al nivel \textit{end-to-end}. Playwright valida que la
interfaz muestra los elementos esperados y reacciona a la actividad del backend,
pero no juzga por sí solo la calidad visual del \textit{dashboard} ni la
claridad de la experiencia de usuario. Esa parte se revisó manualmente durante
el desarrollo.
